\section{Механика твёрдого тела}

Физическое моделирование в данной работе основано на предположении о \emph{абсолютной твёрдости} тел: расстояния между любыми двумя точками тела не изменяются со временем.
Это допущение справедливо для большинства робототехнических задач, где деформации конструкции пренебрежимо малы по сравнению с перемещениями как целого.

\subsection{Состояние твёрдого тела}

Положение твёрдого тела в пространстве полностью описывается шестью обобщёнными координатами:
тремя компонентами вектора положения центра масс $\mathbf{r} \in \mathbb{R}^3$
и тремя параметрами ориентации.
Соответственно, скоростное состояние характеризуется линейной скоростью $\mathbf{v} \in \mathbb{R}^3$
и угловой скоростью $\boldsymbol{\omega} \in \mathbb{R}^3$.
Вектор угловой скорости направлен вдоль мгновенной оси вращения,
его модуль равен угловой скорости вращения.

Для численной симуляции удобно хранить состояние тела в виде вектора:
\begin{equation}
  \mathbf{s} = \bigl(\mathbf{r},\; \mathbf{q},\; \mathbf{v},\; \boldsymbol{\omega}\bigr),
  \label{eq:rigid-state}
\end{equation}
где $\mathbf{q}$ — единичный кватернион, представляющий ориентацию тела.

\subsection{Представление ориентации кватернионами}

Ориентация твёрдого тела определяется матрицей поворота $\mathbf{R} \in SO(3)$, однако хранить $3 \times 3 = 9$ чисел и поддерживать ортогональность матрицы в ходе численного интегрирования нерационально.
На практике применяют единичные кватернионы~\cite{diebel2006representing}:
\begin{equation}
  \mathbf{q} = w + x\,\mathbf{i} + y\,\mathbf{j} + z\,\mathbf{k}, \quad w^2 + x^2 + y^2 + z^2 = 1.
\end{equation}

Кватернион $\mathbf{q}$ представляет поворот на угол $\theta$ вокруг единичной оси $\hat{\mathbf{n}}$:
\begin{equation}
  w = \cos\frac{\theta}{2}, \quad (x,\, y,\, z) = \hat{\mathbf{n}}\sin\frac{\theta}{2}.
\end{equation}

Поворот вектора $\mathbf{p}$ выполняется сэндвич-произведением:
$\mathbf{p}' = \mathbf{q} \otimes \mathbf{p} \otimes \mathbf{q}^{-1}$,
где $\mathbf{p}$ рассматривается как чисто мнимый кватернион $(0, \mathbf{p})$.

Преимущества кватернионов перед углами Эйлера: отсутствие сингулярности «шарнирного замка» (gimbal lock), компактное представление (4 числа), простое интерполирование (SLERP).
По сравнению с матрицами поворота: вдвое меньший объём памяти, возможность простой нормализации для борьбы с численным дрейфом.

\subsection{Уравнения движения}

Динамика твёрдого тела описывается уравнениями Ньютона–Эйлера~\cite{featherstone2008rigid}.

\paragraph{Поступательное движение.}
Второй закон Ньютона для центра масс:
\begin{equation}
  m\dot{\mathbf{v}} = \mathbf{F},
  \label{eq:newton}
\end{equation}
где $m$ — масса тела, $\mathbf{F}$ — суммарная внешняя сила (включая гравитацию и силы контакта).

\paragraph{Вращательное движение.}
В связанной с телом системе координат уравнение Эйлера принимает вид:
\begin{equation}
  \mathbf{I}_{\text{loc}}\dot{\boldsymbol{\omega}} = \boldsymbol{\tau} - \boldsymbol{\omega} \times \bigl(\mathbf{I}_{\text{loc}}\boldsymbol{\omega}\bigr),
  \label{eq:euler-rot}
\end{equation}
где $\mathbf{I}_{\text{loc}} \in \mathbb{R}^{3 \times 3}$ — тензор инерции в связанной системе координат,
$\boldsymbol{\tau}$ — суммарный момент сил.

Для удобства вычислений угловую скорость и момент выражают в мировой системе координат.
Тензор инерции в мировой СК пересчитывается через матрицу поворота:
\begin{equation}
  \mathbf{I} = \mathbf{R}\,\mathbf{I}_{\text{loc}}\,\mathbf{R}^T.
\end{equation}

Тогда уравнение Эйлера в мировой СК:
\begin{equation}
  \dot{\boldsymbol{\omega}} = \mathbf{I}^{-1}\!\left(\boldsymbol{\tau} - \boldsymbol{\omega} \times \mathbf{I}\boldsymbol{\omega}\right).
  \label{eq:euler-world}
\end{equation}

\paragraph{Кинематика.}
Связь скоростей и производных обобщённых координат:
\begin{align}
  \dot{\mathbf{r}} &= \mathbf{v}, \label{eq:kin-r}\\
  \dot{\mathbf{q}} &= \tfrac{1}{2}\,\boldsymbol{\Omega}(\boldsymbol{\omega})\,\mathbf{q}, \label{eq:kin-q}
\end{align}
где $\boldsymbol{\Omega}(\boldsymbol{\omega})$ — кватернионная матрица, соответствующая умножению на $(0,\boldsymbol{\omega})$ слева:
\begin{equation}
  \boldsymbol{\Omega}(\boldsymbol{\omega}) =
  \begin{pmatrix}
    0 & -\omega_x & -\omega_y & -\omega_z \\
    \omega_x & 0 & \omega_z & -\omega_y \\
    \omega_y & -\omega_z & 0 & \omega_x \\
    \omega_z & \omega_y & -\omega_x & 0
  \end{pmatrix}.
\end{equation}

Уравнения~\eqref{eq:newton}–\eqref{eq:kin-q} образуют систему обыкновенных дифференциальных уравнений первого порядка относительно вектора состояния~\eqref{eq:rigid-state}.

\subsection{Тензор инерции}

Тензор инерции характеризует распределение массы тела и определяет связь между угловой скоростью и кинетической энергией вращения:
\begin{equation}
  T_{\text{rot}} = \tfrac{1}{2}\,\boldsymbol{\omega}^T \mathbf{I}\boldsymbol{\omega}.
\end{equation}

Для тела с функцией плотности $\rho(\mathbf{x})$:
\begin{equation}
  I_{ij} = \int_V \rho(\mathbf{x})\left(\|\mathbf{x}\|^2 \delta_{ij} - x_i x_j\right) dV.
  \label{eq:inertia-tensor}
\end{equation}

Диагональные элементы $I_{xx}$, $I_{yy}$, $I_{zz}$ называются \emph{осевыми моментами инерции},
внедиагональные — \emph{центробежными моментами}.
В главных осях инерции тензор диагонален.

На практике тензор инерции вычисляется аналитически для примитивных форм (сфера, коробка, цилиндр) или численно для произвольных сеток.
Хранится в связанной системе координат в диагональной форме ($I_x$, $I_y$, $I_z$), что позволяет работать с тремя числами вместо девяти.

\subsection{Составные тела}

Для тела, составленного из $n$ примитивов с массами $m_k$ и центрами масс $\mathbf{c}_k$:
\begin{equation}
  m = \sum_{k=1}^n m_k, \qquad
  \mathbf{c} = \frac{1}{m}\sum_{k=1}^n m_k \mathbf{c}_k.
\end{equation}

Тензор инерции составного тела вычисляется по теореме Штейнера:
\begin{equation}
  \mathbf{I} = \sum_{k=1}^n \left(\mathbf{I}_k + m_k\left(\|\mathbf{d}_k\|^2\mathbf{E} - \mathbf{d}_k\mathbf{d}_k^T\right)\right),
\end{equation}
где $\mathbf{d}_k = \mathbf{c}_k - \mathbf{c}$ — смещение центра масс $k$-го примитива от общего центра масс, $\mathbf{E}$ — единичная матрица.

\subsection{Выводы по разделу}

Уравнения Ньютона–Эйлера~\eqref{eq:newton},~\eqref{eq:euler-world} совместно с кинематическими уравнениями~\eqref{eq:kin-r},~\eqref{eq:kin-q} образуют полную систему ОДУ, описывающую эволюцию состояния твёрдого тела во времени.
Представление ориентации кватернионами выбрано из-за отсутствия сингулярностей, компактности и удобства нормализации.
В разделе~2.3 рассмотрены методы численного интегрирования этой системы.
